\section{Introduction}
\label{sec:introduction}

Educational decision systems support the allocation of tutoring, counseling, assessments and other scarce resources. Their institutional consequences extend beyond predictive accuracy: students and families need to understand the rule applied to a decision, identify errors and present evidence for reconsideration. An audit therefore needs to connect a case to the policy, inputs and constraints that produced its outcome.

Public controversies illustrate the consequences of weak procedural safeguards. The UK A-level grading algorithm in 2020 generated socially salient downgrades and a policy reversal \cite{Wired2020AlevelAlgorithm}. Wisconsin's dropout early-warning system drew scrutiny over inequitable risk assignments and opaque governance \cite{Markup2023DEWS}. In adjacent public-sector risk profiling, the Dutch SyRI ruling highlighted transparency and accountability requirements \cite{SyRI2020Case}. These cases motivate decision evidence that institutions and affected people can inspect and challenge.

Model-level metrics provide only part of that evidence. Selective indicators can support claims of ethical compliance while leaving responsibility for trade-offs and access to recourse unclear \cite{CitronPasquale2014ScoredSociety,Wieringa2020AccountFor,Metcalf2021AIA,Selbst2019Abstraction}. Educational labels and measurements can also encode structural inequities \cite{BakerHawn2022AlgorithmicBiasEducation}. Governance consequently requires explicit procedures for choosing policies, checking decisions and handling failures.

We present \compa-\ecl-\mas, a layer for translating declared requirements into versioned decision policies and auditable artifacts. ECL checks empirical restrictions; MAS selects a policy from feasible Pareto candidates using ordered criteria and deterministic tie rules. Per-case packets preserve the evidence required to reconstruct the supported decision rule. A verifier checks that evidence, and a technical gate permits export only when verification and the declared empirical constraints pass. Abstention produces a record requiring review; connecting that record to an institutional appeal process remains a deployment responsibility.

\subsection{Contributions}
The paper makes three connected contributions:
\begin{itemize}[leftmargin=*,noitemsep]
\item \textbf{An executable export contract.} We specify how fixed-score replay, evidence verification and empirical feasibility jointly determine technical export permission. Nested acceptance predicates make the conditions and failure reasons explicit.
\item \textbf{Reproducible policy and failure analyses.} Fixed-score sweeps and a retrospective selection/test partition examine policy-dependent feasibility. Seven synthetic fault fixtures test when decision agreement is insufficient for evidence acceptance or export permission.
\item \textbf{Exploratory evidence on structured contestation materials.} A counterbalanced study compares a complete evidence bundle with logs/cards on standardized diagnosis and remediation tasks, with paired estimates and explicit methodological limits.
\end{itemize}

\begin{table*}[t]
\centering\small
\caption{Claims, evidence and explicit boundaries of the revised manuscript.}
\label{tab:claim_evidence}
\begin{tabularx}{\textwidth}{p{.20\textwidth}XX}
\toprule
Claim & Evidence & Boundary \\
\midrule
Fixed-score replay & State reconciliation and taxonomy replay, Tables~\ref{tab:state-counts} and~\ref{tab:policy-replay} & Does not support changed features or medical exceptions. \\
Empirical selection & Algorithm~\ref{alg:mas}, Tables~\ref{tab:run-sweep} and~\ref{tab:recovered} & Same-sample policy exploration; no population fairness guarantee. \\
Explicit failure & Revised runner gate and integration tests, Section~\ref{sec:design} & Historical violation logging was not operational fail-closed; no human queue connected. \\
Structured evidence benefit & Tables~\ref{tab:human-study-summary} and~\ref{tab:human-study-keycases} & Exploratory N=16, no formal ethics review; bundle effect, not isolated replay effect. \\
Entropy diagnostics & Table~\ref{tab:uncertainty_proxy_validation} and structured group-risk report & Scalar proxy; conditional row intervals, not cluster-robust field validation. \\
\bottomrule
\end{tabularx}
\end{table*}


\subsection{Scope}
The contribution concerns governance infrastructure: explicit rules, deterministic policy selection, preserved evidence and observable rejection of exports that fail the contract. The framework combines established selection, provenance and integrity mechanisms; the experiments evaluate their specified acceptance conditions in a reference application. Human task performance and computational integrity are evaluated separately.

The predictors estimate interaction correctness or student success. They provide testbeds for policy mechanics, rather than validated rules for allocating support to the students at greatest risk. The human study evaluates the structured-evidence bundle without isolating its individual components. Accessible appeal channels, legitimate policy adoption, staffing and remedies remain institutional requirements. The demonstrated scope is therefore contestability-enabling infrastructure; prospective early-warning targeting and educational treatment effects require separate evaluation.
